% TODO make hw stylesheet
\documentclass[10pt]{article}
\usepackage[letterpaper,top=1in]{geometry}
\usepackage{tikz-cd}
\usepackage[dvipsnames,svgnames]{xcolor}
\usepackage[shortlabels]{enumitem}
\usepackage[framemethod=TikZ]{mdframed}
\usepackage{amsmath,amssymb,amsthm}
\usepackage[colorlinks]{hyperref}
\usepackage{multicol}
\usepackage{tabularx}
\usepackage{stmaryrd}

\hypersetup{
	colorlinks=true,
	linkcolor=blue,
	filecolor=magenta,
	urlcolor=magenta,
	pdftitle={MAT 534 HW}
}

\usepackage{mathtools}
\usepackage{thmtools}
\usepackage[textsize=footnotesize]{todonotes}
\usepackage{titling}

\reversemarginpar
\mdfdefinestyle{mdbluebox}{roundcorner=10pt,innerbottommargin=9pt,
	linecolor=blue,backgroundcolor=TealBlue!5,}
\declaretheoremstyle[headfont=\sffamily\bfseries\color{MidnightBlue},
mdframed={style=mdbluebox},]{thmbluebox}

\mdfdefinestyle{mdbloobox}{frametitlefont=\bfseries,innerbottommargin=8pt,
	nobreak=true,backgroundcolor=TealBlue!5,linecolor=blue,}
\declaretheoremstyle[headfont=\bfseries\color{MidnightBlue},
mdframed={style=mdbloobox},headpunct={\\[3pt]},postheadspace=0pt,]{thmbloobox}

\mdfdefinestyle{mdredbox}{frametitlefont=\bfseries,innerbottommargin=8pt,
	nobreak=true,backgroundcolor=Salmon!5,linecolor=RawSienna,}
\declaretheoremstyle[headfont=\bfseries\color{RawSienna},
mdframed={style=mdredbox},headpunct={\\[3pt]},postheadspace=0pt,]{thmredbox}

\mdfdefinestyle{mdgreenbox}{linecolor=ForestGreen,backgroundcolor=ForestGreen!5,
	linewidth=2pt,rightline=false,leftline=true,topline=false,bottomline=false,}
\declaretheoremstyle[headfont=\bfseries\sffamily\color{ForestGreen!70!black},
mdframed={style=mdgreenbox},headpunct={ --- },]{thmgreenbox}

\mdfdefinestyle{mdorangebox}{linecolor=RedOrange,backgroundcolor=RedOrange!5,
	linewidth=2pt,rightline=false,leftline=true,topline=false,bottomline=false,}
\declaretheoremstyle[headfont=\bfseries\sffamily\color{RedOrange!70!black},
mdframed={style=mdorangebox},headpunct={ --- },]{thmorangebox}

\mdfdefinestyle{mdblackbox}{linecolor=black,backgroundcolor=RedViolet!5!gray!5,
	linewidth=3pt,rightline=false,leftline=true,topline=false,bottomline=false,}
\declaretheoremstyle[mdframed={style=mdblackbox}]{thmblackbox}

\declaretheoremstyle[spaceabove=6pt,spacebelow=6pt]{basehead}

\declaretheorem[style=basehead,name=Problem]{problem}
\declaretheorem[style=thmredbox,name=Problem,numbered=no]{problem*}
\declaretheorem[style=thmbloobox,name=Setup]{setup} % for config geo
\declaretheorem[style=thmredbox,name=Example,sibling=problem]{example}
\declaretheorem[style=thmredbox,name=$\bigstar$ Problem,sibling=problem]{niceprob}
\declaretheorem[style=thmbluebox,name=Theorem,sibling=problem]{theorem}
\declaretheorem[style=thmbluebox,name=Corollary,sibling=theorem]{corollary}
\declaretheorem[style=thmbluebox,name=Proposition,sibling=theorem]{prop}
\declaretheorem[style=thmbluebox,name=Lemma,numberwithin=problem]{lemma}
\declaretheorem[style=thmbluebox,name=Theorem,numbered=no]{theorem*}
\declaretheorem[style=thmbluebox,name=Lemma,numbered=no]{lemma*}
\declaretheorem[style=thmgreenbox,name=Claim,numberwithin=problem]{claim}
\declaretheorem[style=thmgreenbox,name=Claim,numbered=no]{claim*}
\declaretheorem[style=thmblackbox,name=Remark,numberwithin=problem]{remark}
\declaretheorem[style=thmblackbox,name=Remark,numbered=no]{remark*}
\declaretheorem[style=thmgreenbox,name=Definition,sibling=theorem]{definition}
\declaretheorem[style=thmgreenbox,name=Definition,numbered=no]{definition*}

\providecommand{\ol}{\overline}
\providecommand{\quiv}{\equiv}
\providecommand{\ceil}[1]{\left\lceil #1 \right\rceil}
\providecommand{\floor}[1]{\left\lfloor #1 \right\rfloor}
\newcommand{\dd}{\mathrm{d}}
\providecommand{\ul}{\underline}
\providecommand{\eps}{\varepsilon}
\providecommand{\half}{\frac{1}{2}}
\providecommand{\CC}{\mathbb C}
\providecommand{\FF}{\mathbb F}
\providecommand{\NN}{\mathbb N}
\providecommand{\QQ}{\mathbb Q}
\providecommand{\RR}{\mathbb R}
\providecommand{\ZZ}{\mathbb Z}
\providecommand{\dg}{^\circ}
\providecommand{\ii}{\item}
\providecommand{\setdiff}{\smallsetminus}
\providecommand{\alert}[1]{{\sffamily\textbf{\textcolor{blue}{#1}}}}
\providecommand{\inv}{^{-1}}
\providecommand{\mb}[1]{\mathbf{#1}}

\newenvironment{soln}{\begin{proof}[Solution]}{\end{proof}}

\providecommand{\pow}{\operatorname{Pow}}
\providecommand{\aut}{\operatorname{Aut}}
\providecommand{\ann}{\operatorname{Ann}}
\providecommand{\vocab}[1]{{\textbf{\textcolor{ForestGreen}{#1}}}}
\providecommand{\lcm}{\operatorname{lcm}}
\providecommand{\abs}[1]{\left|#1\right|}
\providecommand{\Ob}{\operatorname{Ob}}
\providecommand{\Hom}{\operatorname{Hom}}
\providecommand{\Aut}{\operatorname{Aut}}
\providecommand{\id}{\operatorname{id}}
\newcommand{\doubleparen}[1]{(\!(#1)\!)}

\usepackage{mathrsfs}

% place title at top right
\setlength{\droptitle}{-8ex}
\pretitle{\begin{flushright}\Large\bfseries}
\posttitle{\par\end{flushright}}
\preauthor{\begin{flushright}}
\postauthor{\end{flushright}}
\predate{\begin{flushright}}
\postdate{\end{flushright}}

\title{MAT 534 HW11}
\author{\large Aditya Pahuja}
\date{\large Due 12/08}
\begin{document}
\maketitle

\begin{problem}
    Let $R$ be a unique factorization domain and $d\in R$ a nonzero element.
    Show that there are only finitely many distinct principal ideals containing $d$.
\end{problem}

\begin{soln}
    Let $u\cdot p_1^{e_1}\cdots p_k^{e_k}$ be a factorization of $d$
    where the $p_i$ are irreducible, the $e_k$ are positive integers,
    and $u$ is a unit.
    Then $d\in (c)$ for some $c\in R$ if and only if $c\mid d$,
    which is true if and only if $c = u'\cdot p_1^{e'_1}\cdots p_k^{e'_k}$
    for nonnegative integers $e'_k\le e_k$ and unit $u'$.
    Two choices of $c$ generate the same ideal if and only if
    the tuples $(e'_1,\dots,e'_k)$ are the same
    (i.e. factorizations are the same up to units),
    and there are $(e_1 + 1)\cdots(e_k+1) < \infty$ choices
    for this tuple ($e_i + 1$ choices for $e'_i\in\{0,1,\dots,e_i\}$),
    hence only finitely many principal ideals containing $d$.
\end{soln}

\begin{problem}
    Let $A\in GL_n(F)$. Prove that there is a polynomial $P(x)\in F[x]$
    of degree at most $n$ such that $A\inv = P(A)$.
\end{problem}

\begin{soln}
    Let $Q(x)\in F[x]$ be the characteristic polynomial of $A$,
    which has degree $n$.
    By the Cayley-Hamilton theorem, $Q(A) = 0$,
    and since $A$ is invertible, $Q(0) = \det A \ne 0$.
    Writing $Q(x) = x^n + a_{n-1}x^{n-1} + \cdots + a_1x + a_0$, define
    \begin{equation*}
	P(x) = -\frac{Q(x) - a_0}{a_0x}
    \end{equation*}
    so that
    \begin{equation*}
	AP(A) = -\frac 1{a_0}Q(A) + I = I
	\iff P(A) = A\inv.\qedhere
    \end{equation*}
\end{soln}

\begin{problem}
    Describe the field of fractions of the ring $\CC[x,y]/(y^2-x^3)$.
\end{problem}

\begin{soln}
    First, observe that the surjective homomorphism
    $\varphi\colon\CC[x,y]\to\CC[t^2, t^3]$
    given by $\varphi(x) = t^2$ and $\varphi(y) = t^3$
    has kernel $(y^2 - x^3)$. Clearly the kernel contains this ideal;
    suppose that $p(x,y)\in\ker\varphi$.
    In $\CC[x,y]/(y^2 - x^3)$, $y^2 = x^3$, so for any polynomial in $\CC[x,y]$,
    we can find a polynomial in the same coset whose $y$-degree is
    reduced to zero or $1$; in $\CC[x,y]$ this corresponds to writing
    $p(x,y) = q(x,y)(y^2 - x^3) + r_1(x)y + r_0(x)$. Then
    \begin{equation*}
        \varphi(p) = r_1(t^2)t^3 + r_0(t^2).
    \end{equation*}
    The degrees of the monomials in $r_1(t^2)t^3$ are all odd,
    while the degrees of the monomials in $r_0(t^2)$ are all even,
    so the only way that the polynomial on the right is zero
    is if $r_1$ and $r_0$ are the zero polynomial, implying
    $p(x,y) = q(x,y)(y^2 - x^3)\in (y^2 - x^3)$.

    Therefore we want to describe the fraction field of
    $\CC[t^2,t^3]$. Since $t = t^3/t^2 = \phi(y)/\phi(x)$
    is in this fraction field, $\operatorname{Frac}(\CC[t^2,t^3])
    = \operatorname{Frac}(\CC[t]) = \CC(t)$.
    In terms of the original ring, the fraction field of $\CC[x,y]/(y^2 - x^3)$
    is $\CC(\ol y/\ol x)$ where $\ol x,\ol y\in\CC[x,y]/(y^2 - x^3)$
    are the projections of $x$ and $y$ to the quotient.
\end{soln}

\pagebreak

\begin{problem}
    Let $T\colon\RR^4\to\RR^4$ be a linear map such that
    the minimal polynomial of $T$ is $x^4 + 1$.
    Determine the number of linear subspaces
    $W\subseteq\RR^4$ such that $T(W)\subseteq W$.
\end{problem}

\begin{soln}
    Treat $\RR^4$ as an $\RR[x]$-module where the action of $x$
    is given by applying $T$. By the classification of
    finitely generated modules over PIDs, we have module isomorphism
    \begin{equation*}
	\RR^4\cong\RR[x]/(x^2 - x\sqrt 2 + 1)\oplus\RR[x]/(x^2 + x\sqrt 2 + 1)
    \end{equation*}
    using that the annihilator of $\RR^4$ is
    $x^4 + 1 = (x^2 - x\sqrt 2 + 1)(x^2 + x\sqrt 2 + 1)$
    and matching the dimensions to choose the exponents of the quadratics.

    Since $T$ has no real eigenvalues ($x^4+1$ has no real roots),
    $T$ does not fix any one-dimensional subspace.

    Among two-dimensional subspaces, there are three cases:
    a two-dimensional subspace of the decomposed module
    is one of $\RR[x]/(x^2 - x\sqrt 2 + 1)$,
    $\RR[x]/(x^2 + x\sqrt 2 + 1)$, or
    a direct sum of a one-dimensional subspace of each factor.
    Since neither one-dimensional subspace is fixed by multiplication by $x$,
    a subspace of the third kind cannot be mapped into itself by $T$,
    while the other two subspaces of course are mapped to themselves
    via multiplication by $x$, giving two $2$-dimensional subspaces
    that are fixed by $T$.

    A $3$-dimensional subspace corresponds to
    a direct sum of one of the factors of the decomposition
    with a one-dimensional subspace of the other factor,
    which cannot be fixed by multiplication by $x$
    because the one-dimensional subspace won't be mapped into itself.

    There are also two trivial subspaces satisfying $T(W)\subseteq W$,
    namely $W = \{0\}$ and $W = \RR^4$, for a total of
    $\boxed 4$ such subspaces $W$.
\end{soln}

\begin{problem}
    Show that a group of order $72$ cannot be simple.
\end{problem}

\begin{soln}
    The number of Sylow $3$-subgroups (of order $9$)
    is $1$ mod $3$ and a divisor of $8$,
    hence either $1$ or $4$. If there is only one then it must be normal,
    so suppose there are four. Then, if $N$ is a normalizer
    of one of these Sylow $3$-subgroups, $\abs G$ does not divide
    $(G:N)! = 24$, which implies that $G$ is not simple by HW02 \#6(d).
\end{soln}

\begin{problem}
    Find, up to isomorphism, all finite groups that have
    precisely three conjugacy classes.
\end{problem}

\begin{soln}
    Let $G$ be a group of order $n\in\NN$ with three conjugacy classes.
    One of these conjugacy classes is $\{1\}$;
    denote the other two by $C_1$ and $C_2$.
    We know that $\abs{C_1}$ and $\abs{C_2}$ are proper divisors of $n$
    and $1 + \abs{C_1} + \abs{C_2} = n$;
    this means that either $1 = \abs{C_1} = \abs{C_2}$
    or one of $1$, $\abs{C_1}$, $\abs{C_2}$ must be larger than $n/3$.

    The former case of course corresponds to $n = 3$, i.e.
    $G\cong\ZZ/3\ZZ$. In the latter case, the only possible proper divisor of $n$
    larger than $n/3$ is $n/2$, if it happens to be an integer;
    this implies that $n$ is even and that
    the larger of $\abs{C_1}$ and $\abs{C_2}$ is $n/2$.
    Supposing without loss of generality that $\abs{C_1} = n/2$,
    we find that $\abs{C_2} = n/2 - 1$, so $n/2 - 1$ also divides $n$.
    Since $\gcd(n/2, n/2 - 1) = 1$, $n$ is divisible by their product
    $n(n-2)/4$, so in particular is at least large as this product. Thus
    \begin{equation*}
	n\ge \frac{n(n-2)}4 \iff 4\ge n-2 \iff n\le 6,
    \end{equation*}
    and we can check that $n=4$ and $n=6$ both are divisible by $n/2 - 1$.
    However, $n=4$ cannot have three conjugacy classes
    because all groups of order four are abelian (either $\ZZ/4\ZZ$ or
    $(\ZZ/2\ZZ)\times(\ZZ/2\ZZ)$), hence have four conjugacy classes.
    Meanwhile, the only non-abelian group of order $6$ is $S_3$
    (easy to check via Sylow theorems),
    which indeed has three conjugacy classes
    \begin{equation*}
	\{(1)\},\quad \{(1\; 2), (1\;3), (2\;3)\},\quad
	\{(1\;2\;3), (1\;3\;2)\}.
    \end{equation*}

    Thus, any finite group with three conjugacy classes
    is isomorphic to either $\boxed{\ZZ/3\ZZ}$ or $\boxed{S_3}$.
\end{soln}

\begin{problem}
    Let $K$ be a field of characteristic $p$
    and $f(x) = x^p - x - a$ for some $a\in K$.
    Show that $f$ is either irreducible in $K[x]$
    or decomposes into a product of linear factors.
\end{problem}

\begin{soln}
    The main idea is the identity
    \begin{equation*}
        f(r + k) = (r+k)^p - (r+k) - a
	= r^k - r - a + k^p - k = r^k - r - a = f(r)
	\tag{$*$}
    \end{equation*}
    for any $k\in\FF^p$ due to Fermat's little theore
    (strictly speaking, $k$ is in the natural embedding of $\FF_p$ in $K$),
    so $r$ is a root if and only if $r+k$ is a root.

    Let $L/K$ be an extension in which $f$ splits
    and let $\alpha$ be one of its roots.
    and let $m(x)\in K[x]$ be the minimal polynomial of $\alpha$ in $K$,
    so the minimal polynomial of $\alpha + k$ is $m(x - k)$.
    If $R_k$ is the set of roots of $m(x-k)$ for $k\in\FF_p$,
    then the union of the $R_k$ is $\{\alpha + k : k\in\FF_p\}$.
    Moreover if two of the $R_k$ intersect then they are the same,
    by the minimality of $m$ (else take the $\gcd$ of the corresponding $m(x-k)$s
    to get a lower degree minimal polynomial for one of the $\alpha + k$).
    In particular the distinct sets among the $R_k$s
    form a partition of the $p$-element set of roots of $f$ in $L$
    into sets of equal size, so the $R_k$s all either have size $1$ or size $p$,
    corresponding respectively to splitting into linear factors
    or being irreducible in $K$.
\end{soln}

\begin{problem}
    Consider the set $\QQ$ of rational numbers as a group under addition.
    Prove that there is no proper subgroup $G\le\QQ$ of finite index.
\end{problem}

\begin{soln}
    Suppose $G\le\QQ$ has finite index $n$
    and let $\pi\colon\QQ\to\QQ/G$ be the projection homomorphism.
    Let $x\in\QQ$ be arbitrary.
    By Lagrange's theorem, $\pi(x/n)$ has order dividing $n$ in $\QQ/G$,
    so $\pi(x) = n\pi(x/n)$ is the identity in $\QQ/G$.
    Thus $x\in\ker\pi = G$, with the equality following from
    the first isomorphism theorem, implying $G$ contains all the rational numbers.
\end{soln}

\begin{problem}
    Let $p$ be a prime number. The symmetric group $S_p$
    acts in a natural way on the set $X = \{1, 2, \dots, p\}$.
    Prove that the induced action of a subgroup $G\le S_p$ is transitive
    if and only if $G$ contains a $p$-cycle.
\end{problem}

\begin{soln}
    If $G$ contains a $p$-cycle then let $\sigma$ be the cycle;
    for any $i,j\in X$, the set $\{i, \sigma(i), \dots, \sigma^{p-1}(i)\}$
    is equal to $X$ because it has $p$ distinct elements
    (since the order of $\sigma$ is $p$), hence
    some power of $\sigma$ yields $j$ when applied to $i$.
    This means the element of $S_p$ corresponding to this power of $\sigma$
    sends $i$ to $j$, i.e. the action is transitive.

    If the induced action of $G$ is transitive,
    then the action of $S_p$ only has one orbit of size $p$.
    By the orbit-stabilizer theorem, $p$ divides the order of $G$,
    so by Cauchy's theorem $G$ contains an element of order $p$.
    The only elements of order $p$ in $S_p$ are the $p$-cycles,
    so we are done.
\end{soln}

\begin{problem}
    Classify, up to isomorphism, all groups of order $45$.
\end{problem}

\begin{soln}
    Let $G$ be a group of order $45$.
    By the Sylow theorems, since $45 = 3^2\cdot 5$,
    the number of subgroups of order $5$ is a divisor of $9$
    and is $1\bmod 5$, which means that $G$ has one subgroup $P_5$ of order $5$.
    Therefore, for any $g$, the subgroup $gP_5g\inv$ also has order $5$,
    so $gP_5g\inv = P_5$, i.e. $P_5$ is normal in $G$.
    Similarly, the number of subgroups of order $9$ is a divisor of $5$
    and is $1\bmod3$, so $G$ has one subgroup $P_9$ of order $9$,
    which by the same logic is normal in $G$.

    We can then see that $P_5P_9 = \{ab : a\in P_5, b\in P_9\}$
    is a subgroup of $G$ because $P_5$ and $P_9$ are normal in $G$,
    and it has order divisible by both $5$ and $9$ because $P_5$, $P_9$
    are subgroups of $P_5P_9$.
    Thus, $P_5P_9$ is a subgroup of $G$ with order at least $45$,
    which means that it is equal to $G$. Furthermore $P_5\cap P_9 = \{1\}$
    because if $P_9$ contained some element $g\in P_5$ not equal to $1$ then,
    since $P_5$ is isomorphic to $\ZZ/5\ZZ$, $g$ generates $P_5$
    so $P_9$ contains $P_5$, which is nonsensical because $5$ doesn't divide $9$.
    To finish we use the following lemma:

    \begin{lemma}
        Suppose $N$ and $K$ are normal in $G$ such that $NK = G$
	and $N\cap K = \{1\}$. Then elements of $N$ commute with elements of $K$,
	and in particular $G\cong N\times K$.
    \end{lemma}

    By this lemma, we find that $G\cong P_5\times P_9$.
    As mentioned above, $P_5\cong \ZZ/5\ZZ$.
    For $P_9$, there are only two cases: either there is an element of order $9$,
    in which case $P_9 \cong \ZZ/5\ZZ$ and $\boxed{G\cong\ZZ/45\ZZ}$,
    or every non-identity element of $P_9$ has order $3$.
    In the latter case, the class equation implies that $3$
    divides the order of the center $\abs{Z(P_5)}$
    (because all non-fixed points have a conjugacy class consisting of $3$ elements).
    If $Z(P_5)$ has only three elements, then suppose $g$
    isn't in $Z(P_5)$; it clearly commutes with itself and every element of $Z(P_5)$,
    so the centralizer of $P_5$ has order at least $4$, meaning it has order $9$,
    but this contradicts $g\notin Z(P_5)$. Therefore $Z(P_5) = P_5$,
    i.e. $P_5$ is abelian. The only possibility, then, is
    $P_5\cong(\ZZ/3\ZZ)\times(\ZZ/3\ZZ)$, so that
    $\boxed{G\cong(\ZZ/15\ZZ)\times(\ZZ/3\ZZ)}$.
\end{soln}

\begin{problem}
    List, up to isomorphism, all $\FF_2[x]$-modules with exactly $8$ elements.
\end{problem}

\begin{soln}
    Let $M$ be an $\FF_2[x]$-module with $8$ elements.
    Since $M$ is finite and $\FF_2[x]$ isn't, $M$ must be a torsion module,
    so it is isomorphic to a direct sum of modules of the form
    $\FF_2[x]/(p(x)^k)$ for some irreducible polynomial $p$
    by the characterization of finitely generated modules over PIDs.
    Noting that $\FF_2[x]/(f(x))$ has $2^{\deg f}$ elements
    because the powers of $x$ in the quotient
    give an $\FF_2$-basis of size $\deg f$,
    the possible decompositions of $M$ are as follows:
    \begin{enumerate}[(1),itemsep=1pt]
	\ii $\FF_2[x]/(x)\oplus\FF_2[x]/(x)\oplus\FF_2[x]/(x)$
	\ii $\FF_2[x]/(x)\oplus\FF_2[x]/(x)\oplus\FF_2[x]/(x+1)$
	\ii $\FF_2[x]/(x)\oplus\FF_2[x]/(x+1)\oplus\FF_2[x]/(x+1)$
	\ii $\FF_2[x]/(x+1)\oplus\FF_2[x]/(x+1)\oplus\FF_2[x]/(x+1)$
	\ii $\FF_2[x]/(x^2)\oplus\FF_2[x]/(x)$
	\ii $\FF_2[x]/(x^2)\oplus\FF_2[x]/(x+1)$
	\ii $\FF_2[x]/((x+1)^2)\oplus\FF_2[x]/(x)$
	\ii $\FF_2[x]/((x+1)^2)\oplus\FF_2[x]/(x+1)$
	\ii $\FF_2[x]/(x^2 + x + 1)\oplus\FF_2[x]/(x)$
	\ii $\FF_2[x]/(x^2 + x + 1)\oplus\FF_2[x]/(x+1)$
	\ii $\FF_2[x]/(x^3 + x + 1)$
	\ii $\FF_2[x]/(x^3 + x^2 + 1)$
	\ii $\FF_2[x]/(x^3)$
	\ii $\FF_2[x]/((x+1)^3)$
    \end{enumerate}
    This list contains all possible decompositions of $M$ because
    the irreducible polynomials in $\FF_2[x]$ of degree $1$ are
    $x$, $x+1$;
    the only irreducible polynomial in $\FF_2[x]$ of degree $2$
    is $x^2 + x + 1$; and
    the irreducible polynomials in $\FF_2[x]$ of degree $3$ are
    $x^3 + x + 1$, $x^3 + x^2 +1$
    (we can verify the latter two fairly quickly by just listing
    out all the polynomials of degree $2$, $3$ and checking for linear factors;
    there are only six polynomials to check because
    the leading and constant coefficients must be $1$).
    Furthermore, any two decompositions in the list are nonisomorphic
    by the structure theorem again, so the list enumerates
    all the isomorphism classes of $8$-elements $\FF_2[x]$-modules.
\end{soln}










\end{document}
